\section{Token, Lexeme, Pattern, dan Atribut}
\label{sec:ch05-token}

\subsection{Tiga Konsep yang Sering Tertukar}

\begin{definisi}[Token, lexeme, dan pattern]
\textbf{Token} adalah kategori leksikal, misalnya \code{IDENT},
\code{INT\_LIT}, atau \code{PLUS}. \textbf{Lexeme} adalah string
konkret yang cocok dengan kategori tersebut, misalnya \code{hasil},
\code{42}, atau \code{+}. \textbf{Pattern} adalah aturan atau
ekspresi reguler yang mendefinisikan bentuk lexeme untuk suatu token.
\end{definisi}

\begin{table}[H]
\centering
\caption{Contoh token, lexeme, dan pattern}
\label{tab:ch05-token-lexeme}
\begin{tabularx}{\textwidth}{@{}l l X@{}}
\toprule
\textbf{Token} & \textbf{Lexeme} & \textbf{Pattern (informal)} \\
\midrule
\code{IDENT} & \code{hasil}, \code{\_x1} &
huruf atau garis bawah, diikuti huruf/angka/garis bawah \\
\code{INT\_LIT} & \code{0}, \code{42} & satu atau lebih digit \\
\code{FLOAT\_LIT} & \code{3.14} & digit, titik, digit \\
\code{PLUS} & \code{+} & karakter \code{+} \\
\code{LE} & \code{<=} & karakter \code{<} diikuti \code{=} \\
\code{STRING\_LIT} & \code{"halo"} & kutip ganda, isi, kutip ganda \\
\bottomrule
\end{tabularx}
\end{table}

Satu token dapat memiliki banyak lexeme. Token \code{IDENT} memiliki
lexeme \code{a}, \code{total}, dan \code{\_tmp}. Sebaliknya, satu
lexeme biasanya dipetakan ke satu token, kecuali kata kunci yang
tampak seperti identifier.

\subsection{Atribut Token}

Token tidak hanya membawa jenis. Ia juga membawa atribut yang
dibutuhkan fase berikutnya:
\begin{itemize}[leftmargin=*]
  \item \textbf{lexeme}: string asli;
  \item \textbf{baris} dan \textbf{kolom}: lokasi untuk diagnostik;
  \item \textbf{nilai}: bilangan bulat, bilangan pecahan, atau kode
        karakter;
  \item \textbf{indeks simbol}: referensi ke tabel simbol bila
        diperlukan.
\end{itemize}

\begin{figure}[H]
\centering
\begin{tikzpicture}
  \node[draw=ubbblue, rounded corners=3pt, fill=ubbblue!5,
        minimum width=11cm, minimum height=2.2cm, align=left,
        font=\small] (tok) {
    \textbf{Token}\\[2pt]
    jenis: \code{TOK\_IDENT} \quad
    lexeme: \code{hasil} \quad
    baris: 4 \quad kolom: 9\\[2pt]
    nilai: --- \quad
    referensi simbol: \code{sym\#17}
  };
\end{tikzpicture}
\caption{Anatomi sebuah token beserta atributnya}
\label{fig:ch05-anatomi}
\end{figure}

\subsection{Kata Kunci dan Identifier}

Kata kunci seperti \code{if}, \code{while}, dan \code{return}
memiliki bentuk yang sama dengan identifier. Ada dua strategi umum.
Pertama, lexer mengenali semua kata sebagai identifier, lalu memeriksa
tabel kata kunci. Kedua, pattern kata kunci ditempatkan sebelum
pattern identifier sehingga aturan kata kunci menang.

Strategi pertama lebih mudah dipelihara karena menambah kata kunci
cukup mengubah tabel. Strategi kedua lebih cepat karena keputusan
diambil langsung oleh DFA. Pada proyek Subset-C, kita memakai
strategi pertama agar tabel kata kunci dapat diuji secara terpisah.

\begin{contoh}[Token dan lexeme]
Untuk potongan program
\begin{lstlisting}[language=C]
int total = 42;
total = total + 1;
\end{lstlisting}
token yang dihasilkan adalah \code{INT}, \code{IDENT(total)},
\code{ASSIGN}, \code{INT\_LIT(42)}, \code{SEMI}, \code{IDENT(total)},
\code{ASSIGN}, \code{IDENT(total)}, \code{PLUS}, \code{INT\_LIT(1)},
\code{SEMI}, \code{EOF}. Lexeme \code{total} muncul tiga kali, tetapi
setiap kemunculan adalah token tersendiri dengan lokasi berbeda.
\end{contoh}
